\section{Benchmark Design}
\label{sec:benchmark}

PrefBench evaluates a seller-facing pricing agent interacting with a buyer over a configurable product bundle. At each episode, the buyer is sampled from a frozen persona bank. The seller observes a limited profile and the current negotiation history, but does not observe the buyer's latent valuation, price sensitivity, patience, or bargaining tendency. The seller then chooses pricing actions over a short negotiation horizon. The episode ends when the buyer accepts, walks away, or the maximum number of rounds is reached.

\subsection{Product and Pricing Substrate}

The current benchmark uses customized vehicle features as the product substrate. Each bundle is composed from a structured catalog, and each feature has both a seller-side implementation cost and a buyer-side contribution to value. The benchmark should not be read as a full model of a real vehicle market. Instead, the catalog provides a controlled setting where feature-level value and pricing decisions can be studied under a fixed contract.

\subsection{Buyer Population}

The buyer population follows a semi-synthetic design. Observable attributes are sampled from configured distributions intended to approximate plausible buyer profiles. Hidden attributes are generated through structured conditional rules that define latent preference and negotiation behavior. This separation is central to the benchmark: agents receive partial buyer information, while the simulator retains hidden variables that affect willingness to pay and responses during negotiation.

\begin{figure}[t]
  \centering
  \includegraphics[width=0.85\linewidth]{ch4/ch4_persona_pipeline.pdf}
  \caption{Persona generation pipeline used by the benchmark prototype. Observable profile variables and hidden buyer traits are generated before episodes are evaluated under a fixed persona bank.}
  \label{fig:persona_pipeline}
\end{figure}

\subsection{Interaction Protocol}

The interaction protocol is shared across agents. Each policy is evaluated on the same test split and seed-controlled episode stream, so differences in outcomes are attributable to policy behavior rather than to different buyer samples. The benchmark records seller actions, buyer responses, termination reason, deal outcome, profit, and round count for each episode.

\begin{figure}[t]
  \centering
  \includegraphics[width=0.85\linewidth]{ch4/ch4_shared_protocol.pdf}
  \caption{Shared benchmark protocol. All agents interact with the same environment contract and are evaluated with the same episode-level metrics.}
  \label{fig:shared_protocol}
\end{figure}

\subsection{Metrics}

The primary metric is average seller profit, because the task is seller-facing pricing. We also report deal rate, walkaway rate, and average rounds. These metrics are reported together because agreement frequency and profit are not equivalent: an agent can close many deals by pricing too softly, while a more conservative agent may preserve margin at the cost of more failed negotiations.
